\diarytitle{0112}{複素指数関数を用いた指数関数と余弦関数の積の不定積分}

オイラーの公式 $e^{i\theta} = \cos\theta + i\sin\theta$ を用いて複素指数関数の積分を計算し、実部・虚部を比較することで実積分の値を求める。

$e^{ax}\cos(bx) = \mathrm{Re}\left(e^{(a+ib)x}\right)$ であるから、まず複素指数関数の積分を計算する。$\lambda = a+ib \neq 0$ として
\[
\int e^{(a+ib)x}\,dx = \frac{1}{a+ib}e^{(a+ib)x} + C_{1}
\]
分母を実数化するため $\dfrac{1}{a+ib} = \dfrac{a-ib}{a^{2}+b^{2}}$ を用いると
\[
\frac{1}{a+ib}e^{(a+ib)x}
= \frac{a-ib}{a^{2}+b^{2}}\,e^{ax}\bigl(\cos bx + i\sin bx\bigr)
\]
右辺を展開して実部を取ると
\[
\mathrm{Re}\left[\frac{a-ib}{a^{2}+b^{2}}\,e^{ax}(\cos bx + i\sin bx)\right]
= \frac{e^{ax}}{a^{2}+b^{2}}\bigl(a\cos bx + b\sin bx\bigr)
\]
したがって
\[
\boxed{\; \int e^{ax}\cos(bx)\,dx = \frac{e^{ax}\bigl(a\cos bx + b\sin bx\bigr)}{a^{2}+b^{2}} + C \;}
\]

（検算: 右辺を $x$ で微分すると
\[
\frac{d}{dx}\left[\frac{e^{ax}(a\cos bx + b\sin bx)}{a^{2}+b^{2}}\right]
= \frac{a\,e^{ax}(a\cos bx + b\sin bx) + e^{ax}(-ab\sin bx + b^{2}\cos bx)}{a^{2}+b^{2}}
\]
分子を整理すると $e^{ax}\bigl[(a^{2}+b^{2})\cos bx + (ab - ab)\sin bx\bigr] = (a^{2}+b^{2})e^{ax}\cos bx$ となるので、全体は $e^{ax}\cos bx$ に一致し、確かに元の被積分関数に戻る。）
